\section{Introduction et philosophie}
\label{sec:intro}

\lettrine[lines=2, findent=2pt, nindent=0pt]{L}{e} marché des changes (FX) est simultanément l'un des plus liquides et l'un des plus difficiles à battre systématiquement. Chaque signal individuel est dominé par un bruit énorme, et les \emph{edges} isolés tendent à s'éroder rapidement lorsqu'ils sont répliqués. La réponse classique à ce problème~--- empiler plusieurs signaux indépendants pour réduire la variance du portefeuille agrégé~--- n'est efficace que si les signaux sont \textbf{vraiment orthogonaux}. Or, la plupart des stratégies FX publiées sont en réalité très corrélées entre elles parce qu'elles exploitent la même source d'alpha sous-jacente (\emph{carry}, tendance $20/50$~EMA, \emph{mean reversion} RSI).

\begin{thesisbox}
\textbf{L'objectif de la Phase 18 est d'assembler trois moteurs de rendement qui capturent des \emph{sources d'alpha structurellement différentes}}, de sorte que leur combinaison réduise le drawdown sans sacrifier le rendement espéré~--- un authentique portefeuille \emph{risk-first}.
\end{thesisbox}

\subsection{Le problème adressé}

Les itérations précédentes (Phases~14 à~17) butaient systématiquement sur le même trade-off~:
\begin{itemize}
\item \textbf{Phase 14}~: combined $2$ sleeves (MR~+~TS), Sharpe IS~$\approx$ 0.88 mais drawdown $>$ 20\,\% et 2019/2023 négatifs.
\item \textbf{Phase 15}~: ajout d'une allocation \emph{regime-adaptive}, qui \emph{dégrade} le Sharpe (0.51) parce que les régimes ne sont pas suffisamment différenciés sur la fenêtre 2019-2025.
\item \textbf{Phase 16}~: retrait du sleeve \emph{cross-section momentum}, contribution négative persistante.
\item \textbf{Phase 17}~: retrait de USD-CAD du sleeve TS Momentum, Sharpe per-pair lift $0.44\rightarrow 0.57$ mais walk-forward 2023 reste marginalement négatif.
\end{itemize}

La Phase~18 franchit le gate en \textbf{ajoutant un troisième sleeve RSI Daily à corrélation négative avec TS Momentum} ($-0.251$) et quasi-nulle avec MR Macro ($-0.027$). L'architecture qui en résulte est illustrée ci-dessous.

\subsection{Architecture du portefeuille}

\begin{figure}[H]
\centering
\small
\begin{tabular}{ccc}
\colorbox{softBlue}{\parbox{4.5cm}{\centering\color{primary}\sffamily\textbf{MR Macro 80\,\%}\\[0.2em]\footnotesize intraday VWAP\\filtres macro (yield curve + chômage)\\EUR-USD uniquement}} &
\colorbox{softYellow}{\parbox{4.5cm}{\centering\color{accent}\sffamily\textbf{TS Momentum 10\,\%}\\[0.2em]\footnotesize daily EMA 20/50 + RSI(7)\\3 paires (EUR/GBP/JPY)\\\textit{sans USD-CAD}}} &
\colorbox{softGreen}{\parbox{4.5cm}{\centering\color{rsiGreen}\sffamily\textbf{RSI Daily 10\,\%}\\[0.2em]\footnotesize daily RSI(14)\\mean reversion\\4 paires equal-weight}} \\[1.2em]
\multicolumn{3}{c}{$\Downarrow$} \\[0.3em]
\multicolumn{3}{c}{\colorbox{lightGrey}{\parbox{12cm}{\centering\color{primary}\sffamily\textbf{Allocation statique 80 / 10 / 10}\\[0.1em]\footnotesize pas d'adaptation de régime, rebalance quotidien implicite}}} \\[0.6em]
\multicolumn{3}{c}{$\Downarrow$} \\[0.3em]
\multicolumn{3}{c}{\colorbox{lightGrey}{\parbox{12cm}{\centering\color{primary}\sffamily\textbf{Vol targeting global}\\[0.1em]\footnotesize \code{target\_vol = 0.28}, \code{max\_leverage = 12}\\$\max(\mathrm{vol}_{21}, \mathrm{vol}_{63})$ avec \code{.shift(1)} anti look-ahead}}} \\[0.6em]
\multicolumn{3}{c}{$\Downarrow$} \\[0.3em]
\multicolumn{3}{c}{\colorbox{primary}{\parbox{12cm}{\centering\color{white}\sffamily\textbf{Équité finale Phase 18}}}} \\
\end{tabular}
\caption{Architecture à trois sleeves orthogonaux du portefeuille Phase~18. Allocation statique, puis vol targeting global, puis pipeline d'équité finale.}
\end{figure}

\subsection{Pourquoi trois sleeves et pas deux ou quatre ?}

\begin{itemize}
\item \textbf{Deux sleeves ne suffisent pas}~: MR~+~TS laisse 2019 et 2023 négatifs (cf. Phase~17).
\item \textbf{Quatre sleeves introduisent du bruit}~: l'ajout d'un sleeve \emph{cross-section momentum} a été testé en Phase~16 et retiré parce que sa corrélation avec TS Momentum est $>0.3$ sans apporter de \emph{lift} sur le Sharpe du combined.
\item \textbf{Trois est le nombre minimal qui couvre les trois régimes dominants}~: \emph{intraday range} (MR Macro), \emph{trend daily} (TS Momentum) et \emph{mean reversion daily} (RSI Daily).
\end{itemize}

\begin{insightbox}
\textbf{La leçon clé de la Phase 18}~: un sleeve avec un Sharpe standalone faible (0.16~pour RSI Daily) peut être le meilleur diversificateur si sa corrélation avec le reste est \emph{négative} sur les années difficiles. Ce qui compte n'est pas le Sharpe individuel mais la \textbf{contribution marginale au Sharpe du portefeuille}, qui dépend de la corrélation.
\end{insightbox}
